% titel: Blatt 4, A3: Wellenlänge, Frequenz, Wellenzahl in Li2+
% skript: Mortimer Kap. 6 (wasserstoffähnliche Atome, Spektrum)

% Teil I: Werkzeuge, kein Lösungsweg.

Für ein wasserstoffähnliches Ion mit Kernladungszahl $Z$ (hier Li$^{2+}$,
$Z=3$) skalieren alle Energieniveaus wie $Z^2$: Die Linien liegen um den
Faktor $Z^2$ höher als beim Wasserstoff.

\begin{werkzeugkasten}[Spektrum wasserstoffähnlicher Atome]
Energiedifferenz beim Übergang $n_i \to n_f$ ($n_i > n_f$), Frequenz und
Wellenzahl:
\begin{equation}
  \nu = R_\infty\,Z^2\left(\frac{1}{n_f^2} - \frac{1}{n_i^2}\right),
  \qquad
  \tilde\nu = \frac{1}{\lambda} = \frac{\nu}{c},
  \qquad \lambda = \frac{c}{\nu}.
\end{equation}
mit $R_\infty \approx \qty{3.29e15}{s^{-1}}$ (Rydberg-Frequenz),
$c = \qty{2.998e8}{m/s}$. Für $Z=1$ (H) ergeben sich die Serien des
Wasserstoffs; für $Z>1$ liegen die gleichen Übergänge bei $Z^2$-facher
Frequenz.
\textit{Skript:} Mortimer Kap. 6 (wasserstoffähnliche Atome,
Formelsammlung), Übung 4/3.
\end{werkzeugkasten}

\begin{achtung}[Darauf achten]
\begin{itemize}
  \item $Z = 3$ nicht vergessen: $\nu \propto Z^2$, die Linie von Li$^{2+}$
        liegt also 9-fach höher (energiereicher) als beim H.
  \item $n_f^2$ und $n_i^2$ nicht vertauschen; benutze $1/n_f^2 - 1/n_i^2$
        (positiv für Emission $n_i > n_f$).
  \item Einheiten: $\nu$ in s$^{-1}$ (Hz), $\tilde\nu$ in m$^{-1}$
        (oft zusätzlich in cm$^{-1}$ umgerechnet), $\lambda$ in m.
  \item Wellenzahl und Frequenz nicht verwechseln: $\tilde\nu = \nu/c$.
\end{itemize}
\end{achtung}
